\documentclass{amsart}
% For dealing with references we use the comment environment
\usepackage{verbatim}
\newenvironment{reference}{\comment}{\endcomment}
%\newenvironment{reference}{}{}
\newenvironment{slogan}{\comment}{\endcomment}
\newenvironment{history}{\comment}{\endcomment}

% For commutative diagrams we use Xy-pic
\usepackage[all]{xy}

% We use 2cell for 2-commutative diagrams.
\xyoption{2cell}
\UseAllTwocells

% We use multicol for the list of chapters between chapters
\usepackage{multicol}

% This is generally recommended for better output
\usepackage{lmodern}
\usepackage[T1]{fontenc}

% For cross-file-references

% Package for hypertext links:
\usepackage{hyperref}

% For any local file, say "hello.tex" you want to link to please

% Theorem environments.
%
\theoremstyle{plain}
\newtheorem{theorem}[subsection]{Theorem}
\newtheorem{proposition}[subsection]{Proposition}
\newtheorem{lemma}[subsection]{Lemma}

\theoremstyle{definition}
\newtheorem{definition}[subsection]{Definition}
\newtheorem{example}[subsection]{Example}
\newtheorem{exercise}[subsection]{Exercise}
\newtheorem{situation}[subsection]{Situation}

\theoremstyle{remark}
\newtheorem{remark}[subsection]{Remark}
\newtheorem{remarks}[subsection]{Remarks}

\numberwithin{equation}{subsection}

% Macros
%
\def\lim{\mathop{\mathrm{lim}}\nolimits}
\def\colim{\mathop{\mathrm{colim}}\nolimits}
\def\Spec{\mathop{\mathrm{Spec}}}
\def\Hom{\mathop{\mathrm{Hom}}\nolimits}
\def\Ext{\mathop{\mathrm{Ext}}\nolimits}
\def\SheafHom{\mathop{\mathcal{H}\!\mathit{om}}\nolimits}
\def\SheafExt{\mathop{\mathcal{E}\!\mathit{xt}}\nolimits}
\def\Sch{\mathit{Sch}}
\def\Mor{\mathop{\mathrm{Mor}}\nolimits}
\def\Ob{\mathop{\mathrm{Ob}}\nolimits}
\def\Sh{\mathop{\mathit{Sh}}\nolimits}
\def\NL{\mathop{N\!L}\nolimits}
\def\CH{\mathop{\mathrm{CH}}\nolimits}
\def\proetale{{pro\text{-}\acute{e}tale}}
\def\etale{{\acute{e}tale}}
\def\QCoh{\mathit{QCoh}}
\def\Ker{\mathop{\mathrm{Ker}}}
\def\Im{\mathop{\mathrm{Im}}}
\def\Coker{\mathop{\mathrm{Coker}}}
\def\Coim{\mathop{\mathrm{Coim}}}

% Boxtimes
%
\DeclareMathSymbol{\boxtimes}{\mathbin}{AMSa}{"02}

%
% Macros for moduli stacks/spaces
%
\def\QCohstack{\mathcal{QC}\!\mathit{oh}}
\def\Cohstack{\mathcal{C}\!\mathit{oh}}
\def\Spacesstack{\mathcal{S}\!\mathit{paces}}
\def\Quotfunctor{\mathrm{Quot}}
\def\Hilbfunctor{\mathrm{Hilb}}
\def\Curvesstack{\mathcal{C}\!\mathit{urves}}
\def\Polarizedstack{\mathcal{P}\!\mathit{olarized}}
\def\Complexesstack{\mathcal{C}\!\mathit{omplexes}}
% \Pic is the operator that assigns to X its picard group, usage \Pic(X)
% \Picardstack_{X/B} denotes the Picard stack of X over B
% \Picardfunctor_{X/B} denotes the Picard functor of X over B
\def\Pic{\mathop{\mathrm{Pic}}\nolimits}
\def\Picardstack{\mathcal{P}\!\mathit{ic}}
\def\Picardfunctor{\mathrm{Pic}}
\def\Deformationcategory{\mathcal{D}\!\mathit{ef}}

\setlength{\emergencystretch}{3em}
\begin{document}
\title[Stacks Project, Tag 0BCA]{353 --- Lifting a prescribed projective complex across an ideal}
\maketitle
\noindent Copyright (C) 2005--2025 Johan de Jong. Stacks Project authors. Source: \href{https://stacks.math.columbia.edu/tag/0BCA}{Tag 0BCA}.

\noindent Editorial difficulty note (not author text). Difficulty H2: The author lifts an acyclic mapping cone and then constructs a comparison map whose degreewise cokernels stay in the chosen projective class. A separate split-injectivity argument removes the remaining obstruction, producing a complex that simultaneously represents the original derived object and has the prescribed reduction..

\noindent Editorial provenance note (not author text). History: excerpt from revision \texttt{a0444\allowbreak 6e57e\allowbreak c1fbc\allowbreak 252a8\allowbreak 71afc\allowbreak ec775\allowbreak 2fb28\allowbreak 07b14}, more-algebra.tex, lines 19389--19466. Reference presentation was adapted; original-author context and core support blocks were retained or added. The mathematical wording is unchanged. Licensed under GNU FDL 1.2 or later; no invariant sections or cover texts. See ../licenses/COPYING. Context blocks reproduce author definitions and supporting results; they are not additional counted problems.

\begin{definition}
\label{derived-definition-cone}
Let $\mathcal{A}$ be an additive category.
Let $f : K^\bullet \to L^\bullet$ be a morphism of
complexes of $\mathcal{A}$. The {\it cone} of $f$
is the complex $C(f)^\bullet$ given by
$C(f)^n = L^n \oplus K^{n + 1}$ and
differential
$$
d_{C(f)}^n =
\left(
\begin{matrix}
d^n_L & f^{n + 1} \\
0 & -d_K^{n + 1}
\end{matrix}
\right)
$$
It comes equipped with canonical morphisms of complexes
$i : L^\bullet \to C(f)^\bullet$ and $p : C(f)^\bullet \to K^\bullet[1]$
induced by the obvious maps $L^n \to C(f)^n \to K^{n + 1}$.
\end{definition}

\begin{lemma}
\label{lemma-lift-complex}
Let $R$ be a ring. Let $I \subset R$ be an ideal. Let $\mathcal{P}$
be a class of $R$-modules.
Let $K \in D(R)$ and let $E^\bullet$ be a complex of $R/I$-modules
representing $K \otimes_R^\mathbf{L} R/I$. Assume
\begin{enumerate}
\item each $P \in \mathcal{P}$ is a projective $R$-module,
\item $P_1 \in \mathcal{P}$ and $P_1 \oplus P_2 \in \mathcal{P}$
if and only if $P_1, P_2 \in \mathcal{P}$,
\item if $f : P_1 \to P_2$, $P_1, P_2 \in \mathcal{P}$ is surjective
modulo $I$, then $f$ is surjective,
\item $E^\bullet$ is bounded above and $E^i$ is of the form $P/IP$
for $P \in \mathcal{P}$, and
\item $K$ can be represented by a bounded above complex
whose terms are in $\mathcal{P}$.
\end{enumerate}
Then there exists a bounded above complex $P^\bullet$ whose terms
are in $\mathcal{P}$ with $P^\bullet/IP^\bullet$ isomorphic to
$E^\bullet$ and representing $K$ in $D(R)$.
\end{lemma}

\begin{proof}
By assumption (5) we can represent $K$ by a bounded above complex 
$K^\bullet$ whose terms are in $\mathcal{P}$. Then $K \otimes_R^\mathbf{L} R/I$
is represented by $K^\bullet/IK^\bullet$. Since $E^\bullet$ is
a bounded above complex of projective $R/I$-modules by (4), 
we can choose a quasi-isomorphism $\delta : E^\bullet \to K^\bullet/IK^\bullet$
(Derived Categories, Lemma
\href{https://stacks.math.columbia.edu/tag/064B}{lemma morphisms from projective complex (Tag 064B)}).
Let $C^\bullet$ be cone on $\delta$
(Derived Categories, Definition \ref{derived-definition-cone}).
The module $C^i$ is the direct sum $K^i/IK^i \oplus E^{i + 1}$
hence is of the form $P/IP$ for some $P \in \mathcal{P}$
as (2) says in particular that $\mathcal{P}$ is preserved under taking sums.
Since $C^\bullet$ is acyclic, we can apply
Lemma \href{https://stacks.math.columbia.edu/tag/0BC9}{lemma lift acyclic complex (Tag 0BC9)} and find a acyclic
lift $A^\bullet$ of $C^\bullet$. The complex $A^\bullet$ is
bounded above and has terms in $\mathcal{P}$. In
$$
\xymatrix{
K^\bullet \ar@{..>}[r] \ar[d] & A^\bullet \ar[d] \\
K^\bullet/IK^\bullet \ar[r] & C^\bullet \ar[r] & E^\bullet[1]
}
$$
we can find the dotted arrow making the diagram commute
by Derived Categories, Lemma
\href{https://stacks.math.columbia.edu/tag/0649}{lemma morphisms lift projective (Tag 0649)}.
We will show below that it follows from (1), (2), (3)
that $K^i \to A^i$ is the inclusion of a direct summand
for every $i$. By property (2) we see that $P^i = \Coker(K^i \to A^i)$
is in $\mathcal{P}$. Thus we can take
$P^\bullet = \Coker(K^\bullet \to A^\bullet)[-1]$ to conclude.

\medskip\noindent
To finish the proof we have to show the following: Let $f : P_1 \to P_2$,
$P_1, P_2 \in \mathcal{P}$ and $P_1/IP_1 \to P_2/IP_2$ is split
injective with cokernel of the form $P_3/IP_3$ for some
$P_3 \in \mathcal{P}$, then $f$ is split injective.
Write $E_i = P_i/IP_i$. Then $E_2 = E_1 \oplus E_3$.
Since $P_2$ is projective we can choose a map $g : P_2 \to P_3$
lifting the map $E_2 \to E_3$. By condition (3) the map $g$
is surjective, hence split as $P_3$ is projective. Set $P_1' = \Ker(g)$
and choose a splitting $P_2 = P'_1 \oplus P_3$. Then $P'_1 \in \mathcal{P}$
by (2). We do not know that
$g \circ f = 0$, but we can consider the map
$$
P_1 \xrightarrow{f} P_2 \xrightarrow{projection} P'_1
$$
The composition modulo $I$ is an isomorphism. Since $P'_1$ is
projective we can split $P_1 = T \oplus P'_1$. If $T = 0$, then
we are done, because then $P_2 \to P'_1$ is a splitting of $f$.
We see that $T \in \mathcal{P}$ by (2).
Calculating modulo $I$ we see that $T/IT = 0$.
Since $0 \in \mathcal{P}$ (as the summand of any $P$ in $\mathcal{P}$)
we see the map $0 \to T$ is surjective and we conclude that $T = 0$
as desired.
\end{proof}
\end{document}
